\section{Related Work}

\noindent\textbf{Cybercrime and Fraud Detection.}
Early cybercrime detection relied on feature-engineered models for phishing and fraud 
identification~\cite{cybercrime2020}. While effective in constrained settings, these 
approaches struggle with scalability and generalization across diverse linguistic and 
contextual variations prevalent in India.

\noindent\textbf{Multilingual Transformer Models.}
BERT~\cite{bert2019} and its multilingual variants mBERT and XLM-R~\cite{xlmr2020} 
enabled strong cross-lingual transfer across low-resource languages, making them 
well-suited for linguistically diverse settings. However, existing cybercrime detection 
works adopting these models remain text-only, limiting applicability to real-world 
incidents that frequently involve screenshots and fraudulent visual content.

\noindent\textbf{Multimodal Learning.}
CLIP~\cite{clip2021} and Vision Transformers~\cite{vit2021} established strong 
foundations for joint visual-textual representation. Subsequent models including 
Flamingo~\cite{flamingo2022}, BLIP-2~\cite{blip2_2023}, LLaVA~\cite{llava2023}, 
Qwen2-VL~\cite{qwen2vl2024}, and SigLIP~\cite{siglip2023} further advanced 
cross-modal alignment through instruction tuning. Baltru\v{s}aitis et al.~\cite{multimodal2019} 
survey multimodal fusion strategies comprehensively. While such frameworks have been 
applied to misinformation detection~\cite{fake_news_2025}, their use in multilingual 
cybercrime detection remains limited.

\noindent\textbf{Explainable AI.}
LIME~\cite{lime2016} provides local model-agnostic interpretability, while 
DIME~\cite{dime2022} and VALE~\cite{vale2024} extend explanations to multimodal 
settings. Explanation-guided training is explored in~\cite{megl2024}. Surveys 
by~\cite{xai_transformer2025, vit_xai2024, vlm_survey2024} confirm that despite 
rapid XAI progress, its integration into multimodal cybercrime pipelines remains 
largely unexplored.

\noindent\textbf{Research Gaps.}
Three critical gaps motivate this work: (i) no large-scale multilingual multimodal 
cybercrime dataset exists for the Indian context; (ii) existing detection systems 
cover a narrow set of crime categories; and (iii) no prior work unifies multilingual 
understanding, multimodal fusion, and explainability in a single framework. This work 
directly addresses all three through a 25,000-sample dataset, a 50-category taxonomy, 
and an end-to-end transformer-based architecture with integrated XAI.